\section{Nonvanishing on a prime-power sequence}\label{sec:nonvanishing}

Fix \(M\in4\mathbb N\), set \(h=M/2\), and choose an odd prime \(p>2M+h\).
For every \(k\ge0\) put \(q=p^k\), \(m=Mq\), \(g=hq\). The scaled pencil
satisfies
\begin{equation}\label{eq:nv-pencil}
\begin{aligned}
 \mathcal H_k(X)&:=q^{3m}G_m(X/q^2)
 =\det\begin{bmatrix}qC^{(g)}\\q^2B^{(g)}-XA^{(g)}\end{bmatrix}
 \in\mathbb Z_{(p)}[X],\\
 \overline{\mathcal H_k}(X)&=\overline{G_M}(X)^q.
\end{aligned}
\end{equation}
Every nonzero denominator in the formulas of Section~\ref{sec:source} has
absolute value less than \((2M+h)q<pq\). So \(q/d\) is \(p\)-integral. It
reduces to zero unless \(q\mid d\), and to \(e^{-1}\) at \(d=eq\),
\(0<|e|<p\). Likewise \(q^2/(d_1d_2)\) vanishes on reduction unless both
denominators are divisible by \(q\).

Write \(i=aq+v\), \(j=bq+w\), \(0\le v,w<q\). Lucas and Frobenius give
\begin{equation}\label{eq:nv-lucas}
 \binom{(a+h)q+v}{hq}\equiv\binom{a+h}{h},\qquad
 U_{aq+v}^{(hq)}\equiv2^v U_a^{(h)}\pmod p.
\end{equation}
Multiply by \(q\). In the first sum in \(t_i\), only the denominator
indices \(dq\) with \(1\le d\le a\) survive. In the second, only \(dq\) with
\(1\le d\le h\) survive. Since \(q\) is odd, parity is preserved. Also
\(2^{dq}-1\equiv2^d-1\pmod p\), and Lucas reduces the binomial to the seed
one. Therefore \(qt_{aq+v}^{(hq)}\equiv2^vt_a^{(h)}\pmod p\). If \(v\ne
w\), then \(q\nmid j-i\), so both \(qC_{ij}\) and \(q^2B_{ij}\) reduce to
zero. If \(v=w\) and \(a\ne b\), the divided differences reduce to \(2^v\)
times their seed entries.

On the diagonal, the multiples of \(q\) in \([i+1,i+hq]\) are exactly
\((a+1)q,\ldots,(a+h)q\). This holds even when \(v=0\) or \(q-1\). Thus
\(qC_{ii}\equiv2^vC_{aa}^{(h)}\). In the first diagonal \(B\) sum only the
squared denominators \((dq)^2\) survive. In the double sum both indices must
be multiples of \(q\), with \(2\le d\le h\), \(1\le e<d\). Lucas, parity and
Frobenius reduce these terms to the seed double sum. Hence
\(q^2B_{ii}\equiv2^vB_{aa}^{(h)}\) and \(A_{ii}\equiv2^vA_{aa}^{(h)}\). All
scaled entries are integral before reduction. The argument holds for each
\(k\) directly, without a recurrence in \(k\).

The window endpoints are multiples of \(q\). Grouping rows and columns by
residue gives, in each residue block, the seed pencil times \(2^v\). The
row and column permutations are the same transpose permutation of a
\(2M\times q\) array, so their signs cancel. The product of the scalar
block determinants is \(2^{Mq(q-1)}\equiv1\pmod p\), because
\(p-1\mid q-1\). This proves \eqref{eq:nv-pencil}, including the exact
scalar and sign.

The modular data in the supplementary materials show that for \(M=32\),
\(p=83\) the matrix of Section~\ref{sec:source} satisfies
\begin{equation}\label{eq:nv-seed}
 \det\begin{bmatrix}C^{(16)}\\B^{(16)}\end{bmatrix}
                           \equiv30\pmod{83}.
\end{equation}
This is a finite \(64\times64\) matrix check. The saved entries, row swaps
and unnormalized Gaussian pivots determine the determinant by integer
modular arithmetic. The code also rebuilds the matrix with exact rationals,
and independently by a modular Pascal recurrence. No larger determinant is
computed.

Let \(M=32,p=83\), and fix \(t\in\mathbb Q_{83}\) with
\(2k+v_{83}(t)>0\). Then \(X=q^2t\) lies in \(83\mathbb Z_{83}\). By
\eqref{eq:nv-pencil} and \eqref{eq:nv-seed},
\begin{equation}\label{eq:nv-final}
 q^{3m}G_m(t)=\mathcal H_k(q^2t)\equiv30^q=30\pmod{83},
 \qquad v_{83}(G_{32\cdot83^k}(t))=-3k(32\cdot83^k).
\end{equation}
This proves \eqref{eq:nonvanishing}. For \(t=a/b\) in lowest terms,
\(k\ge\lfloor v_{83}(b)/2\rfloor+1\) suffices; \(v_p(0)=+\infty\) covers
zero. The sequence does not depend on \(t\), and every index lies in
\(32\mathbb N\).
